\documentclass[../ap_2014.tex]{subfiles}

\graphicspath{{./image/}}

\begin{document}

\setcounter{chapter}{1}
\chapter{微分方程式}
\section{}
\subsection{}
定数係数斉次線形微分方程式なので解を\(y=Ae^{\lambda t}\)とおいて、\(\lambda\)を求める。
\begin{equation}\begin{aligned}[b]
    \qty(\dv[3]{t}-3\dv[2]{t}+2\dv{t})Ae^{\lambda t} &= 0\\
    \qty(\lambda^3-3\lambda^2+2\lambda)Ae^{\lambda t} &= 0\\
    \lambda(\lambda-1)(\lambda-2) &= 0\\
    \lambda &= 0, 1, 2
\end{aligned}\end{equation}
よって求める一般解はこれらのときの重ね合わせである
\begin{equation}\begin{aligned}[b]
    y = C_1 + C_2e^{t}+C_3e^{2t}
\end{aligned}\end{equation}
である。

\subsection{}
前問の斉次微分方程式が非斉次になったものなので、
特解を加えたものが求める一般解である。
特解として\(y = at^3+bt^2+ct\)を仮定すると
\begin{equation}\begin{aligned}[b]
    3t^2 &= \qty(\dv[3]{t}-3\dv[2]{t}+2\dv{t})\qty(at^3+bt^2+ct)\\
    &= 6a-3(6at+2b)+2(3at^2+2bt+c)\\
    &= 6at^2 +(-18a+4b)t+6a-6b+2c
\end{aligned}\end{equation}
これより
\begin{equation}\begin{aligned}[b]
    a &= \frac{1}{2}, & b &= \frac{9}{4}, & c &=\frac{21}{4}
\end{aligned}\end{equation}

以上より一般解は
\begin{equation}\begin{aligned}[b]
    y = C_1 + C_2e^{t}+C_3e^{2t} + \frac{1}{2}t^3+\frac{9}{4}t^2+\frac{21}{4}t
\end{aligned}\end{equation}

\section{}
\subsection{}
(2)式を行列で表すと
\begin{equation}\begin{aligned}[b]
    &\begin{cases}
        \dv{x}{t} &= \dot{x}\\
        \dv{\dot{x}}{t} &= -2\dot{x}-3y+2\\
        \dv{y}{t} &= \dot{x}+2y
    \end{cases}\\
    &\rightarrow
    \dv{t}\begin{pmatrix}
        x(t)\\
        \dot{x}(t)\\
        y(t)
    \end{pmatrix}
    =\underset{=A}{\underline{\begin{pmatrix}
        0 & 1 & 0\\
        0 & -2 & -3\\
        0 & 1 & 2
    \end{pmatrix}}}
    \begin{pmatrix}
        x(t)\\
        \dot{x}(t)\\
        y(t)
    \end{pmatrix}
    +\underset{=b}{\underline{\begin{pmatrix}
        0\\
        2\\
        0
    \end{pmatrix}}}
\end{aligned}\end{equation}
また。初期条件は
\begin{equation}\begin{aligned}[b]
    c = (0\,0\,2)^\mathsf{T}
\end{aligned}\end{equation}
のように表される。

\subsection{}
\(A\)の固有値方程式より固有値は
\begin{equation}\begin{aligned}[b]
    0 &= \abs{A-\lambda} = \lambda^3-\lambda\\
    \lambda &= 1, 0, -1
\end{aligned}\end{equation}
である。

\(\lambda_0=0\)の固有ベクトル\(v_0\)は
\begin{equation}\begin{aligned}[b]
    \begin{pmatrix}
        0 & 1 & 0\\
        0 & -2 & -3\\
        0 & 1 & 2
    \end{pmatrix}\begin{pmatrix}
        x\\
        y\\
        z
    \end{pmatrix}=0
\end{aligned}\end{equation}
より\(v_0 = (1\,0\,0)^\mathsf{T}\)である。

\(\lambda_+=1\)の固有ベクトル\(v_+\)は
\begin{equation}\begin{aligned}[b]
    \begin{pmatrix}
        -1 & 1 & 0\\
        0 & -3 & -3\\
        0 & 1 & 1
    \end{pmatrix}\begin{pmatrix}
        x\\
        y\\
        z
    \end{pmatrix}=0
\end{aligned}\end{equation}
より\(v_+ = (1\,1\,-1)^\mathsf{T}\)である。

\(\lambda_-=-1\)の固有ベクトル\(v_-\)は
\begin{equation}\begin{aligned}[b]
    \begin{pmatrix}
        1 & 1 & 0\\
        0 & -1 & -3\\
        0 & 1 & 3
    \end{pmatrix}\begin{pmatrix}
        x\\
        y\\
        z
    \end{pmatrix}=0
\end{aligned}\end{equation}
より\(v_- = (0\,3\,-1)^\mathsf{T}\)である。

これより\(A\)を対角化する行列は
\begin{equation}\begin{aligned}[b]
    P = \begin{pmatrix}
        1 & 1 & 0\\
        0 & 1 & 3\\
        0 & -1 & -1
    \end{pmatrix}
\end{aligned}\end{equation}
となる。
これを使うと
\begin{equation}\begin{aligned}[b]
    P^{-1}AP &= \text{diag}(0,1,-1)\\
    A &= P\text{diag}(0,1,-1)P^{-1}
\end{aligned}\end{equation}
より
\begin{equation}\begin{aligned}[b]
    e^{tA} &= e^{tP\text{diag}(0,1,-1)P^{-1}}=P\text{diag}(1,e^{t},e^{-t})P^{-1}\\
    &=\begin{pmatrix}
        1 & 1 & 0\\
        0 & 1 & 3\\
        0 & -1 & -1
    \end{pmatrix}
    \begin{pmatrix}
        1 & 0 & 0\\
        0 & e^{t} & 0\\
        0 & 0 & e^{-t}
    \end{pmatrix}
    \begin{pmatrix}
        1 & 1/2 & 3/2 \\
        0 & -1/2 & -3/2\\
        0 &1/2 & 1/2
    \end{pmatrix}\\
    &=\begin{pmatrix}
        1 & 1 & 0\\
        0 & 1 & 3\\
        0 & -1 & -1
    \end{pmatrix}
    \begin{pmatrix}
        1 & 1/2 & 3/2\\
        0 & -e^{t}/2 & -3e^{t}/2\\
        0 & e^{-t}/2 &e^{-t}/2
    \end{pmatrix}\\
    &= \begin{pmatrix}
        1 & (1-e^{t})/2 & 3(1-e^{t})/2\\
        0 & -e^{t}/2+3e^{-t}/2 & -3e^{t}/2+3e^{-t}/2\\
        0 & e^{t}/2-e^{-t}/2 & 3e^{t}/2-e^{-t}/2
    \end{pmatrix}
\end{aligned}\end{equation}

\subsection{}
問題1の線形微分方程式の斉次解は\(C_0\)を任意のベクトルとして
\begin{equation}\begin{aligned}[b]
    x(t) = e^{tA}C_0
\end{aligned}\end{equation}
となる。
また、特殊解の形を
\begin{equation}\begin{aligned}[b]
    x(t) = \begin{pmatrix}
        pt + q\\
        rt + s\\
        ut + v
    \end{pmatrix}
\end{aligned}\end{equation}
とおくと微分方程式の両辺は
\begin{equation}\begin{aligned}[b]
    \dv{x}{t} &= Ax+b\\
    \begin{pmatrix}
        p\\
        r\\
        u
    \end{pmatrix}
    &= \begin{pmatrix}
        0 & 1 & 0\\
        0 & -2 & -3\\
        0 & 1 & 2
    \end{pmatrix}
    \begin{pmatrix}
        pt + q\\
        rt + s\\
        ut + v
    \end{pmatrix}
    +\begin{pmatrix}
        0 \\
        2 \\
        0
    \end{pmatrix}\\
    &= \begin{pmatrix}
        rt+s\\
        (-2r-u)t-2s-3v+2\\
        (r+2u)t +s+2v
    \end{pmatrix}
\end{aligned}\end{equation}
これより
\begin{equation}\begin{aligned}[b]
    x(t)=\begin{pmatrix}
        4t+q\\
        4\\
        -2
    \end{pmatrix}
\end{aligned}\end{equation}
よって微分方程式の解は
\begin{equation}\begin{aligned}[b]
    x(t) = e^{tA}C_0
    +\begin{pmatrix}
        4t+q\\
        4\\
        -2
    \end{pmatrix}
\end{aligned}\end{equation}
\(x(0)=c\)でなければならないので
\begin{equation}\begin{aligned}[b]
    x(0)&= c\\
    \begin{pmatrix}
        C_{0}^{(1)}+q\\
        C_{0}^{(2)}+4\\
        C_{0}^{(3)}-2
    \end{pmatrix}
    &=\begin{pmatrix}
        0 \\
        0 \\
        2
    \end{pmatrix}
\end{aligned}\end{equation}
以上より
\begin{equation}\begin{aligned}[b]
    x(t) &= \begin{pmatrix}
        1 & (1-e^{t})/2 & 3(1-e^{t})/2\\
        0 & -e^{t}/2+3e^{-t}/2 & -3e^{t}/2+3e^{-t}/2\\
        0 & e^{t}/2-e^{-t}/2 & 3e^{t}/2-e^{-t}/2
    \end{pmatrix}\begin{pmatrix}
        -q\\
        -4\\
        4
    \end{pmatrix}
    +\begin{pmatrix}
        4t+q\\
        4\\
        -2
    \end{pmatrix}
    =\begin{pmatrix}
        4(1+t-e^{t})\\
        4(1-e^{t})\\
        4e^{t}-2
    \end{pmatrix}
\end{aligned}\end{equation}

\end{document}
